% Part: lambda-calculus
% Chapter: introduction
% Section: lambda-computable

\documentclass[../../../include/open-logic-section]{subfiles}

\begin{document}

\olfileid{lam}{int}{cmp}
\olsection{الدوال \usetoken{S}{lambda definable} قابلة للحساب}

\begin{thm}
\ollabel{thm:lambda-computable}
إذا كانت دالة جزئية~$f$ !!{lambda defined} بحد لامبدا، فهي قابلة للحساب.
\end{thm}

\begin{proof}
افترض أن دالة~$f$ !!{lambda defined} بحد لامبدا~$X$. ولنصف إجراءً غير
رسمي لحساب~$f$. عند إدخال~$m_0$، \dots،~$m_{n-1}$، اكتب الحد
$X \num m_0 \ldots \num m_{n-1}$. ابن شجرة بكتابة جميع اختزالات الحد
الأصلي ذات الخطوة الواحدة أولًا؛ واكتب تحتها جميع الاختزالات ذات الخطوة
الواحدة لتلك الحدود (أي اختزالات الحد الأصلي ذات الخطوتين)؛ واستمر على
هذا النحو. فإذا بلغت عددًا، فأعده جوابًا؛ وإلا كانت الدالة غير معرّفة.

ويخبرنا الاحتكام إلى أطروحة تشيرش أن هذه الدالة قابلة للحساب. والطريقة
الأفضل لإثبات المبرهنة هي إعطاء وصف عودي لإجراء البحث هذا. فيمكن، مثلًا،
تعريف متتالية من الدوال والعلاقات العودية البدائية:
``$\fn{IsASubterm}$'' و``$\fn{Substitute}$''
و``$\fn{ReducesToInOneStep}$'' و``$\fn{ReductionSequence}$''
و``$\fn{Numeral}$''، وما إلى ذلك. وعندئذ يكون الإجراء العودي الجزئي
لحساب~$f(m_0, \dots, m_{n-1})$ هو البحث عن متتالية من الاختزالات ذات
الخطوة الواحدة تبدأ بـ$X \num{m_0} \dots \num{m_{n-1}}$ وتنتهي بعدد،
ثم إعادة العدد الموافق لذلك العدد. والتفاصيل طويلة ومضجرة، لكنها روتينية.
\end{proof}

\end{document}
